% Converted from S1_ThyClinScore_sensitivity/reviewer_response_package_20261120_claude_v24/RESPONSE_R1_5_EN.md (2026-11-17)
We thank the reviewer for these suggestions. We have added a blinded clinician assessment, sensitivity analyses for the lesion-matching threshold and all score weights, and an audit of the template libraries against the test sets.\par

\textbf{Blinded clinician assessment.} Two clinicians independently rated 270 reports of 90 cases (30 per test set; the ThyroidXAgent report and two baseline reports per case), blinded to the system, against the reference report on five dimensions (clinical accuracy, factual consistency, completeness, readability and clinical usability; 1--5). Inter-rater agreement (ICC(2,1)) across the three test sets was 0.54--0.78 for clinical accuracy, 0.60--0.78 for factual consistency, 0.57--0.79 for clinical usability, 0.33--0.57 for completeness and 0.08--0.25 for readability (Supplementary Table~\SuppTabClinVal{}). The Spearman correlation of ThyClinScore with the clinical-accuracy rating was 0.437 (95\% CI 0.271--0.586) in SMU-HMC, 0.464 (0.234--0.663) in ZJH-TS and 0.560 (0.373--0.709) in KMVE.\par

\textbf{Sensitivity to the matching threshold and weights.} We recomputed ThyClinScore for all models on the three test sets under 120 parameter configurations (\responsefigureref{fig_rep_sensitivity_curves} and \responsefigureref{fig_rep_sensitivity_ranks}; KMVE results use the revised KMVE library described below). These varied the lesion-matching threshold (0.3--0.7), the weight of size relative to feature agreement in lesion matching ($\alpha_s$ 0.5--0.9), the size tolerances for gland and lesion measurements ($\tau$ 0.10--0.30 and 0.15--0.45), $\lambda$ (0.1--0.5), $\eta$ (0.3--0.7) and each of the six consistency weights and five completeness weights (0.5--1.5 times the default), singly and in three joint grids ($\lambda$ with $\eta$, the threshold with $\alpha_s$, and the two tolerances). ThyroidXAgent had the highest mean ThyClinScore in all 120 configurations in SMU-HMC, in 110 in KMVE and in 117 in ZJH-TS, including every variation of the component weights in all three test sets. In KMVE, LLaVA-Med, which reported no abnormality in any KMVE case, ranked first in the ten configurations with a matching threshold of 0.6 or 0.7 (differences $-$0.049 to $-$0.029; 95\% CIs excluding zero). In ZJH-TS, ThyroidXAgent ranked second or third in three configurations, and all of these differences had 95\% CIs that included zero. The ranking of all models changed little (median Kendall's $\tau_b$ 0.94 in SMU-HMC, 0.93 in KMVE and 0.94 in ZJH-TS), except in KMVE when the matching threshold was 0.3 or size and feature agreement were given equal or nearly equal weight in lesion matching (minimum $\tau_b$ 0.21).\par

\textbf{Patient overlap.} The SMU-HMC library was built after the 400 test patients had been excluded, and none of its 7,147 source reports shares a report or patient identifier with the SMU-HMC test set. KMVE follows the official split, in which no training identifier matches a test identifier.\par

\textbf{Near-duplicate report text.} We compared each complete test reference report with every source report of its library, using the Jaccard similarity of character 3-grams after removing section headings, punctuation and spaces, and counted a test report as having a near-duplicate when the similarity reached 0.90. In SMU-HMC, 4 of the 400 test reports had a near-duplicate source report; these four reports and their five matching source reports all describe a normal thyroid without focal lesions in the standard wording used at the centre, and they differ only in the gland measurements. No ZJH-TS test report had a near-duplicate. In KMVE, 133 of the 492 test reports had a near-duplicate training report (86 describing a normal gland, 40 a nodule and 7 a postoperative thyroid), and 105 had normalised text identical to that of a training report, most of them (83) describing a normal gland. We used the training and test split published with KMVE without modification. The frequent identical texts are probably because the released KMVE reports mask sizes and other measurements and retain only short descriptions of the findings, so that reports of normal glands and of simple nodules are often worded identically across patients. These identical texts led us to revise the KMVE library, as described next.\par

\textbf{Revision of the KMVE library.} We excluded the 406 training reports that were near-duplicates of KMVE test reports and rebuilt the KMVE library with the procedure used for the SMU-HMC library (our response to Reviewer 2, Comment \hyperref[rev:R2.19]{19}), using the remaining 1,313 training reports to build the templates and the 246 validation reports to check them before the test set was scored. On the KMVE test set (\responsetableref{tab_rep_kmve_revision}), ROUGE-L increased from 0.5880 to 0.6573 (difference 0.0694; 95\% CI 0.0597 to 0.0791), BLEU-1 to BLEU-4 decreased by 0.035 to 0.068, METEOR decreased by 0.017, and ThyClinScore changed from 0.4407 to 0.4374 ($-$0.0033; 95\% CI $-$0.0209 to +0.0144). ThyroidXAgent remained first among the methods of Supplementary Tables~\SuppTabNLG{} and \SuppTabClinical{} on all six lexical metrics and on ThyClinScore. The KMVE results in the manuscript have been updated accordingly.\par

\textbf{Removal of near-duplicate templates.} From the SMU-HMC library, which is used for both SMU-HMC and ZJH-TS reports, we removed every template with a similarity of at least 0.90 to any sentence of at least 20 characters in the near-duplicate test reports identified above (5 templates) and regenerated all reports from the same tool outputs. None of these templates had been used in an SMU-HMC or ZJH-TS report, so these reports did not change. In a stricter analysis, we applied the same criterion to every test report and removed 166 templates. Lexical metrics were computed as in Supplementary Table~\SuppTabNLG{} and ThyClinScore as in Supplementary Table~\SuppTabClinical{} (\responsetableref{tab_rep_template_removal}). ROUGE-L changed by $-$0.0016 (95\% CI $-$0.0031 to $-$0.0004) in SMU-HMC and +0.0003 ($-$0.0015 to +0.0022) in ZJH-TS, and ThyClinScore by $-$0.0034 ($-$0.0067 to $-$0.0005) and +0.0043 ($-$0.0005 to +0.0106), respectively. In both analyses, ThyroidXAgent kept its rank among the methods of Supplementary Tables~\SuppTabNLG{} and \SuppTabClinical{} on all six lexical metrics and on ThyClinScore in both test sets.\par

\textbf{Centre-specific template effects.} To assess the effect of the template wording, we exchanged the clause templates between the SMU-HMC library and the revised KMVE library while keeping the tool outputs and the report format of each test set (Supplementary Table~\SuppTabSwap{} and \responsefigureref{fig_rep_template_exchange}). KMVE reports were generated with the SMU-HMC library and, as in KMVE reports, kept only the gland, nodule and lymph-node sentences, with nodule sizes masked. SMU-HMC and ZJH-TS reports used gland and nodule sentences generated with the KMVE library and filled with the measured sizes, together with their original measurement sentence, lymph-node sentence and impression. The lexical metrics decreased substantially: ROUGE-L by 0.1247 (95\% CI 0.1175 to 0.1318) in SMU-HMC, 0.2895 (0.2775 to 0.3017) in KMVE and 0.1045 (0.0936 to 0.1155) in ZJH-TS, and BLEU-4 by 0.137 to 0.276. ThyClinScore decreased much less: by 0.0484 (0.0409 to 0.0558) in SMU-HMC and 0.0471 (0.0289 to 0.0646) in KMVE, and in ZJH-TS it changed by $-$0.0082 ($-$0.0217 to +0.0055). This indicates that exchanging the templates mainly changed how the findings were worded and organized, while most of the clinically relevant content was retained. The remaining decrease reflects findings that the exchanged templates state differently or omit; for example, the KMVE templates, which follow the KMVE reports, do not describe the aspect ratio or internal echo of a nodule. Lexical metrics count shared word sequences, so they also fall when the same findings are worded or ordered differently. For example, one SMU-HMC nodule was described with the original templates as ``hypoechoic area seen in the right thyroid lobe, about 43.5 $\times$ 13.0 mm, margin clear, shape regular, aspect ratio < 1, internal echo fairly homogeneous; CDFI: no flow signal within it'' and with the KMVE templates as ``a hypoechoic nodule is seen in the right lobe, about 43.5 $\times$ 13.0 mm, with a well-defined margin and a regular shape; CDFI shows no detectable flow signal'' (translated from the Chinese reports). Both sentences state a hypoechoic nodule in the right lobe of 43.5 $\times$ 13.0 mm with a clear margin, a regular shape and no flow signal, but the ROUGE-L of the report fell from 0.4940 to 0.3783 and its BLEU-4 from 0.2172 to 0.1347, whereas its ThyClinScore was 0.5341 and 0.5349. ThyClinScore compares the findings after the report has been structured, which makes it much less sensitive to reporting style; this is why we designed it. All metrics in this audit were computed with one protocol for every condition. We added this analysis to Methods~\SecMethReport{}.\par
